\chapter{Small-capacity punctures and removable stability}\label{sec:puncture-stability}
The cutoffs here shrink around a missing point, in contrast with the large-radius cutoffs used in Bernstein arguments. We work on an already smooth manifold with a locally bounded potential; removal of a singularity of the hypersurface itself is a different regularity problem. The geometric source is \cite{ChMin}, Lemma 33.5, pp.~81--82, and the cutoff discussion in \cite{Sharp17}, p.~324.

\begin{definition}[Local cutoffs avoiding a closed set]\label{def:local-zero-capacity-cutoffs}
A closed measure-zero subset $E$ of a smooth manifold $\Sigma$ has the cutoff property used here if, for every compact $K\subset\Sigma$, there are locally Lipschitz functions $0\leq\eta_j\leq1$ on a neighborhood of $K$, vanishing on an open neighborhood of $E\cap K$, such that $\eta_j\to1$ almost everywhere on $K$ and $\int_K|\nabla\eta_j|^2\to0$. The property must hold also on larger compact neighborhoods containing the support of any test function. This explicit local zero-capacity condition is all that is used below; broader equivalences for capacity are not assumed. Compare the cutoff mechanism in \cite{ChMin}, Lemma 33.5.
\end{definition}

\begin{lemma}[A point has cutoffs of arbitrarily small energy in dimension at least two]\label{lem:point-capacity-cutoffs}
On a smooth $n$-dimensional Riemannian manifold with $n\geq2$, a single interior point has the property in Definition \ref{def:local-zero-capacity-cutoffs}. The cutoffs can equal one outside a shrinking coordinate ball. Their Dirichlet energies are $O(\varepsilon^{n-2})$ for $n>2$ and $O(1/|\log\varepsilon|)$ for $n=2$.
\end{lemma}
\begin{proof}
Use normal coordinates centered at the point, where the metric and volume density are uniformly comparable with their Euclidean counterparts in a fixed small ball. For $n>2$, take a radial cutoff zero for $r\leq\varepsilon$, one for $r\geq2\varepsilon$, and with derivative bounded by $2/\varepsilon$. Its energy is at most $C\varepsilon^{-2}\varepsilon^n$.
For $n=2$ and $0<\varepsilon<1$, take zero for $r\leq\varepsilon^2$, one for $r\geq\varepsilon$, and
\[
 \eta_\varepsilon(r)=\frac{\log(r/\varepsilon^2)}{\log(1/\varepsilon)}
 \quad(\varepsilon^2<r<\varepsilon).
\]
Its energy is bounded by
\[
 \frac{C}{\log(1/\varepsilon)^2}\int_{\varepsilon^2}^{\varepsilon}\frac{dr}{r}
 =\frac{C}{\log(1/\varepsilon)}.
\]
Extend each cutoff by one outside the chart. They are locally Lipschitz, vanish near the point, and converge to one away from it. A point has zero Riemannian volume, so the almost-everywhere requirement follows. The estimates are the small-radius logarithmic cutoff argument of \cite{ChMin}, Lemma 33.5.
\end{proof}

\begin{lemma}[Finite unions and approximation of test functions]\label{lem:capacity-test-approximation}
Finite sets on a smooth manifold of dimension at least two have the cutoff property. More generally, if a closed set $E$ has that property, every $\phi\in C_c^\infty(\Sigma)$ can be approximated in $H^1$ by functions $\phi_j\in C_c^\infty(\Sigma\setminus E)$ whose supports lie in a fixed compact set. For any locally bounded $q$, the quadratic forms $\int(|\nabla\phi_j|^2-q\phi_j^2)$ converge to the form of $\phi$.
\end{lemma}
\begin{proof}
For finitely many points take the product of their cutoffs. Each factor lies in $[0,1]$, so for $N$ factors its squared gradient is at most $N$ times the sum of the squared gradients. Pointwise convergence off the finite set and vanishing energy follow from Lemma \ref{lem:point-capacity-cutoffs}.
For a general $E$, choose a compact neighborhood $K$ of $\operatorname{spt}\phi$ and cutoffs there. The functions $\phi\eta_j$ vanish near $E$ on their support and
\[
 \nabla(\phi\eta_j-\phi)=(\eta_j-1)\nabla\phi+\phi\nabla\eta_j.
\]
Dominated convergence controls the first term and the $L^2$ error, while the vanishing energy controls the second term. Smooth approximation of each compactly supported Lipschitz function inside the open set $\Sigma\setminus E$, with error tending to zero in $H^1$, gives the required smooth sequence. The supports can remain in one fixed compact neighborhood. Since $q$ is bounded there, $H^1$ convergence also gives convergence of the potential term, by $L^2$ convergence and Cauchy--Schwarz.
\end{proof}

\begin{theorem}[Stability extends across sets with the cutoff property]\label{thm:removable-stability-capacity}
Let $\Sigma$ be smooth, $q\in L^\infty_{\mathrm{loc}}(\Sigma)$, and let $E\subset\Sigma$ be closed with the property in Definition \ref{def:local-zero-capacity-cutoffs}. If
\[
 \int_\Sigma(|\nabla\phi|^2-q\phi^2)\geq0
 \quad\text{for all }\phi\in C_c^\infty(\Sigma\setminus E),
\]
then the same inequality holds for all $\phi\in C_c^\infty(\Sigma)$. In particular, on a smooth two-sided minimal hypersurface of dimension at least two, stability away from finitely many points implies stability everywhere. This supplies the smooth-manifold cutoff deduction in \cite{ChMin}, Lemma 33.5, and \cite{Sharp17}, p.~324.
\end{theorem}
\begin{proof}
Approximate an arbitrary compactly supported smooth test by Lemma \ref{lem:capacity-test-approximation}. Each approximant has nonnegative quadratic form, and the forms converge, so the limiting form is nonnegative. For a smooth minimal hypersurface the Jacobi potential $q=|B|^2+\operatorname{Ric}(\nu,\nu)$ is bounded on compact subsets. Finite sets satisfy the cutoff property in dimension at least two, so the general result applies. This proof does not require extending a positive Jacobi function across the missing points, nor does it prove that an initially singular hypersurface becomes smooth.
\end{proof}

\begin{example}[Removing a point can change stability in dimension one]\label{ex:point-removal-one-dimensional}
On $(-1,1)$ the quadratic form
\[
 Q(f)=\int_{-1}^1\left(|f'|^2-\frac{\pi^2}{2}f^2\right)
\]
is nonnegative on $C_c^\infty((-1,1)\setminus\{0\})$ but negative on some $C_c^\infty(-1,1)$ test.
\end{example}
\begin{proof}
On each unit-length half-interval the positive sine function solves $w''+\pi^2w=0$. Applying Proposition \ref{prop:positive-supersolution-transform} with potential $\pi^2$ gives $\int|f'|^2\geq\pi^2\int f^2$ for compactly supported tests there. Splitting a test supported away from zero into its two halves proves nonnegativity for the displayed smaller potential. On the full interval, $f(x)=\cos(\pi x/2)$ belongs to $H_0^1(-1,1)$ and has $\int|f'|^2=(\pi^2/4)\int f^2$, so $Q(f)<0$. Smooth compactly supported approximations in $H_0^1$ retain strict negativity. Thus the dimension restriction in point-capacity removal cannot be discarded.
\end{proof}

\begin{corollary}[Ordered sheets away from finitely many points imply global stability]\label{cor:ordered-sheets-punctured-stability}
Let $\Sigma^n$ be a smooth connected closed two-sided embedded minimal hypersurface, $n\geq2$, and let $P\subset\Sigma$ be finite. Suppose on an exhaustion of $\Sigma\setminus P$ there are two globally ordered minimal normal graphs whose functions converge smoothly to zero on compact subsets. Then $\Sigma$ is stable. In particular this configuration is impossible if the ambient Ricci curvature is strictly positive along $\Sigma$.
\end{corollary}
\begin{proof}
A connected manifold of dimension at least two remains connected after finitely many points are removed: a path meeting those points can be rerouted in small punctured coordinate balls. Theorem \ref{thm:separated-sheets-jacobi} therefore gives a positive Jacobi function on $\Sigma\setminus P$, and Proposition \ref{prop:positive-jacobi-stability} gives stability there. Theorem \ref{thm:removable-stability-capacity} extends it across $P$. If the ambient Ricci curvature is positive, the global test one instead gives
\[
 Q(1)=-\int_\Sigma\bigl(|B|^2+\operatorname{Ric}(\nu,\nu)\bigr)\,dA<0,
\]
contradicting stability. The existence of globally ordered normal sheets and a smooth two-sided limit are hypotheses; numerical multiplicity alone has not established them here.
\end{proof}
